
\cleardoublepage
\phantomsection
\section*{Copyright and terms of use of the work by third parties}

This is an academic work that can be used by third parties as long as internationally accepted rules and good practices are respected with regard to copyright and related rights.
The work may be used under the terms of the licence below.
If a use is not covered by that licence, the user should contact the author through the RepositóriUM of the University of Minho.

\noindent\textbf{Licence granted to users of this work}\\
Attribution-NonCommercial\\
CC BY-NC\\
\url{https://creativecommons.org/licenses/by-nc/4.0/}

\cleardoublepage
\phantomsection
\section*{Acknowledgements}

Looking back on the work that led to this dissertation, I am deeply grateful to the people whose guidance, support and encouragement made it possible.

My deepest thanks go to my supervisor, Professor Doctor Paulo Alexandre Ribeiro Cortez, who has been a constant source of motivation throughout this research.
His knowledge, rigour and dedication to academic excellence guided my work at every stage, and I learned a great deal from his teaching and advice during the master's programme and the whole period of this study.

I am also grateful to my co-supervisor, Professor Doctor Emanuel Demétrio Mendes Gouveia, for his insights and encouragement.
His perspective and feedback were essential in shaping and broadening the scope of this work, and the collaborative atmosphere he fostered contributed greatly to my academic growth.

I also thank the telecommunications operator that made the network data available, without which this study would not have been possible, and the developers of the open-source tools on which this work is built, in particular PonyGE2 and STUMPY.

Finally, I thank my family and friends for their unconditional love, patience and encouragement throughout this journey.

\cleardoublepage
\phantomsection
\section*{Statement of integrity}

I hereby declare having conducted this academic work with integrity.
I confirm that I have not used plagiarism or any form of undue use of information or falsification of results along the process leading to its elaboration.
I further declare that I have fully acknowledged the Code of Ethical Conduct of the University of Minho.

\cleardoublepage
\phantomsection
\section*{Abstract}
\addcontentsline{toc}{section}{Abstract}

This dissertation focuses on anomaly detection in grouped telecommunications time-series data from a real network.
Modern telecommunications networks generate large volumes of traffic that must be monitored to detect anomalous behaviour that can affect service reliability, capacity provisioning, and operational decision-making.
These datasets are structured along several categorical dimensions, which creates many possible aggregation levels and makes it difficult to decide which time series to analyse.
This work proposes a framework based on hierarchical aggregation structures, with the detector held fixed.
A single global view or fully disaggregated monitoring is often inadequate: a global view can hide local anomalies, while full disaggregation is costly and can leave many series too sparse for reliable detection.
Instead, this work explores partial hierarchical structures that may reveal abnormal traffic patterns at a relevant level of detail, aiming to detect anomalies in complex grouped datasets without sacrificing interpretability or incurring prohibitive analytical cost.

A formal grammar automatically generates hierarchical structures.
Each candidate is scored with the Matrix Profile, and the area under the ROC curve (AUC) of those scores is the fitness of the search.
In addition to predefined hierarchies, the method can discover useful aggregation structures automatically.
To this end, Grammatical Evolution (\ac{GE}) \parencite{oneill2001grammatical,oneill2003} is used.
The main contribution is a reproducible, interpretable framework for searching aggregation structures with a fixed detector for telecommunications time series.
The evolved hierarchies reached a test AUC of $0.678 \pm 0.066$, roughly 0.10 to 0.14 above the fully aggregated series and the Call, Event, and Service Type splits ($p = 0.002$).

\noindent\textbf{Keywords:} anomaly detection, telecommunications, hierarchical time series, grouped time series, Grammatical Evolution, Matrix Profile.

\cleardoublepage
\phantomsection
\section*{Resumo}
\addcontentsline{toc}{section}{Resumo}

Esta dissertação foca-se na deteção de anomalias em dados de séries temporais agrupadas de telecomunicações provenientes de uma rede real.
As redes de telecomunicações modernas geram grandes volumes de tráfego que devem ser monitorizados para detetar comportamento anómalo que pode afetar a fiabilidade do serviço, o aprovisionamento de capacidade e a tomada de decisão operacional.
Estes conjuntos de dados estão estruturados ao longo de várias dimensões categóricas, o que cria muitos níveis de agregação possíveis e dificulta a decisão sobre quais séries temporais analisar.
Este trabalho propõe um enquadramento baseado em estruturas de agregação hierárquica, mantendo o detetor fixo.
Uma vista global única ou uma monitorização totalmente desagregada é frequentemente inadequada: uma vista global pode ocultar anomalias locais, enquanto a desagregação total é dispendiosa e pode deixar muitas séries demasiado esparsas para uma deteção fiável.
Em vez disso, este trabalho explora estruturas hierárquicas parciais que podem revelar padrões anómalos de tráfego a um nível de detalhe relevante, com o objetivo de detetar anomalias em conjuntos de dados agrupados complexos sem sacrificar a interpretabilidade nem incorrer num custo analítico proibitivo.

Uma gramática formal gera automaticamente estruturas hierárquicas.
Cada candidata é pontuada com o Matrix Profile, e a área sob a curva ROC (AUC) dessas pontuações é a aptidão da procura.
Para além de hierarquias predefinidas, o método pode descobrir automaticamente estruturas de agregação úteis.
Para este efeito, recorre-se à Evolução Gramatical (\ac{GE}) \parencite{oneill2001grammatical,oneill2003}.
A principal contribuição é um enquadramento reproduzível e interpretável para a procura de estruturas de agregação com um detetor fixo em séries temporais de telecomunicações.
As hierarquias evoluídas atingiram uma AUC de teste de $0{,}678 \pm 0{,}066$, cerca de 0,10 a 0,14 acima da série totalmente agregada e das divisões por Call, Event e Service Type ($p = 0{,}002$).

\noindent\textbf{Palavras-chave:} deteção de anomalias, telecomunicações, séries temporais hierárquicas, séries temporais agrupadas, Evolução Gramatical, Matrix Profile.

\setcounter{tocdepth}{1}
\tableofcontents

\cleardoublepage
\phantomsection
\section*{List of abbreviations and acronyms}
\addcontentsline{toc}{section}{List of abbreviations and acronyms}
\begin{acronym}[]
    \acro{AUC}{\textit{Area Under the ROC Curve}}
    \acro{BNF}{\textit{Backus--Naur Form}}
    \acro{CDR}{\textit{Call Detail Record}}
    \acro{CNN}{\textit{Convolutional Neural Network}}
    \acro{GA}{\textit{Genetic Algorithm}}
    \acro{GE}{\textit{Grammatical Evolution}}
    \acro{GP}{\textit{Genetic Programming}}
    \acro{ISP}{\textit{Internet Service Provider}}
    \acro{KPI}{\textit{Key Performance Indicator}}
    \acro{LSTM}{\textit{Long Short-Term Memory}}
    \acro{MAE}{\textit{Mean Absolute Error}}
    \acro{MP}{\textit{Matrix Profile}}
    \acro{PIG}{\textit{Position-Independent Grow}}
    \acro{ROC}{\textit{Receiver Operating Characteristic}}
    \acro{ROI}{\textit{Region of Interest}}
    \acro{VAE}{\textit{Variational Autoencoder}}
\end{acronym}

\cleardoublepage
\phantomsection
\addcontentsline{toc}{section}{\listfigurename}
\listoffigures

\cleardoublepage
\phantomsection
\addcontentsline{toc}{section}{\listtablename}
\listoftables
